% =====================================================================
% VERIFICA DI MATEMATICA - 2BS
% File A e B: 75 minuti e 30 punti per fila.
% Punteggi per fila: 4 + 4 + 4 + 10 + 6 + 2 = 30.
% Compilare almeno due volte con pdfLaTeX.
% Per mostrare le soluzioni: decommentare \printanswers.
% =====================================================================

\documentclass[11pt,a4paper,addpoints]{exam}

% ---------- Lingua e codifica ----------
\usepackage[utf8]{inputenc}
\usepackage[T1]{fontenc}
\usepackage[italian]{babel}
\usepackage{lmodern}
\usepackage{microtype}

% ---------- Matematica e impaginazione ----------
\usepackage{amsmath,amssymb}
\usepackage{xcolor}
\usepackage{tikz}
\usepackage{array,tabularx}
\usepackage[
  a4paper,
  top=2cm,
  bottom=2.2cm,
  left=2cm,
  right=2.2cm,
  headheight=30pt,
  headsep=12pt
]{geometry}

% ---------- Soluzioni ----------
%\printanswers
\ifdefined\CORRETTORE
  \printanswers
\fi
\cancelspace
\unframedsolutions
\renewcommand{\solutiontitle}{%
  \noindent\textbf{Soluzione.}\enspace
}
\SolutionEmphasis{\color{blue!60!black}}

% ---------- Dati della prova ----------
\newcommand{\materia}{Matematica}
\newcommand{\classe}{2BS}
\newcommand{\indirizzo}{Liceo Scientifico}
\newcommand{\tipoprova}{Verifica di matematica}
\newcommand{\argomento}{Disequazioni e sistemi di disequazioni}
\newcommand{\durata}{75 minuti}
\newcommand{\annoscolastico}{2026/2027}
\newcommand{\docente}{Prof. Rubens Longhi}
\newcommand{\fila}{A}

\newcommand{\indicazioni}{%
  Svolgi sul foglio a protocollo e mostra i passaggi.
  Per ogni disequazione o sistema scrivi l'insieme delle soluzioni.
  Lavora in $\mathbb{R}$, senza calcolatrice.
}

% ---------- Macro matematiche ----------
\newcommand{\R}{\mathbb{R}}

% ---------- Intestazioni e piè di pagina ----------
\pagestyle{headandfoot}
\firstpageheader
  {%
    \scriptsize\textsc{IIS David Maria Turoldo}\\[-1pt]
    \scriptsize Via Ronco 11 -- 24019 Zogno (BG)
  }
  {}
  {%
    \scriptsize A.S.\ \annoscolastico\\[-1pt]
    \scriptsize \docente
  }
\firstpageheadrule
\runningheader
  {\small\materia\ -- \classe}
  {\small\tipoprova}
  {\small Fila \fila}
\runningheadrule
\newcommand{\paginefila}{\pageref{fine@fila@\fila}}
\firstpagefooter{}{\small Pagina \thepage\ di \paginefila}{}
\runningfooter{}{\small Pagina \thepage\ di \paginefila}{}

% ---------- Obiettivi e formato degli esercizi ----------
\newcommand{\obiettivocorrente}{}
\newcommand{\obiettivo}[1]{%
  \gdef\obiettivocorrente{#1}%
}
\newcommand{\stampaobiettivo}{%
  \ifx\obiettivocorrente\empty
  \else
    {\footnotesize\color{black!55}%
      \textit{Obiettivo: \obiettivocorrente}}%
  \fi
}
\qformat{%
  \textbf{\thequestion.}\enspace
  \stampaobiettivo
  \hfill
  {\small(\totalpoints\ \points)}%
  \gdef\obiettivocorrente{}%
}
\renewcommand{\questionlabel}{}
\renewcommand{\partlabel}{\thequestion\thepartno)}
\renewcommand{\subpartlabel}{\thesubpart.}
\pointpoints{punto}{punti}
\renewcommand{\partshook}{%
  \setlength{\itemsep}{2pt}%
  \setlength{\rightmargin}{0.5cm}%
}

% ---------- Numerazione indipendente per ciascuna fila ----------
\newcounter{offsetfila}
\newcounter{totalqfila}
\renewcommand{\thequestion}{%
  \number\numexpr
    \value{question}-\value{offsetfila}%
  \relax
}
\newenvironment{esercizi}
  {%
    \begin{questions}
    \setcounter{question}{\value{offsetfila}}%
  }
  {\end{questions}}

% ---------- Esercizi non spezzati nella versione studenti ----------
\newenvironment{esercizio}
  {%
    \par
    \ifprintanswers\else\begin{samepage}\fi
  }
  {%
    \ifprintanswers\else\end{samepage}\fi
    \par\vspace{3pt}
  }

% ---------- Tabella dei punteggi ----------
\hqword{\textbf{Esercizio}}
\hpword{\textbf{Punti massimi}}
\hsword{\textbf{Punti ottenuti}}
\htword{\textbf{Totale}}

% ---------- Inizio e fine di ciascuna fila ----------
\newcommand{\iniziofila}[1]{%
  \clearpage
  \renewcommand{\fila}{#1}%
  \setcounter{page}{1}%
  \setcounter{offsetfila}{\value{totalqfila}}%
  \begin{center}
    {\small\bfseries
      \tipoprova
      \ifprintanswers
        -- Soluzioni e criteri di correzione
      \fi
    }\\[3pt]
    {\Large \argomento}\\[2pt]
    {\small
      \indirizzo\ -- Classe \classe\
      -- Durata: \durata\
      -- \textbf{Fila #1}
    }
  \end{center}
  \vspace{4pt}
  \noindent
  \begin{tabularx}{\textwidth}{@{}l X l p{3cm}@{}}
    \textbf{Nome e Cognome:} & \hrulefill
    & \textbf{Data:} & \hrulefill\\[8pt]
  \end{tabularx}
  \noindent
  \fbox{%
    \parbox{\dimexpr\textwidth-2\fboxsep-2\fboxrule}{%
      \small
      \textbf{Indicazioni.} \indicazioni\space
      Punteggio totale: \pointsinrange{fila#1} punti.
    }%
  }
  \medskip
  \begingradingrange{fila#1}%
}
\newcommand{\finefila}{%
  \setcounter{totalqfila}{\value{question}}%
  \endgradingrange{fila\fila}%
  \ifprintanswers
  \else
    \vfill
    \begin{center}
      \small
      \partialgradetable{fila\fila}[h][questions]
    \end{center}
    \vspace{6pt}
    \noindent
    {\small
      \textbf{Voto:}\ \rule{3cm}{0.4pt}
      \hfill
      \textbf{Firma del docente:}\ \rule{5cm}{0.4pt}
    }%
  \fi
  \label{fine@fila@\fila}%
}

\begin{document}
\linespread{1}\selectfont
\ifprintanswers\else
\setlength{\abovedisplayskip}{5pt}
\setlength{\belowdisplayskip}{5pt}
\setlength{\abovedisplayshortskip}{3pt}
\setlength{\belowdisplayshortskip}{3pt}
\fi
\ifprintanswers\small\fi
\iniziofila{A}
\begin{esercizi}
\begin{esercizio}
\obiettivo{valutare proprietà e procedimenti}
\question Indica se ogni affermazione è vera o falsa. Giustifica la risposta.
\begin{parts}
\part[1] È possibile risolvere la disequazione $\dfrac{x-1}{x}>\dfrac{2x-1}{x}$ moltiplicando ambo i membri per $x$.
\begin{solution}
Falso. Il segno di $x$ è variabile: se $x<0$ occorre invertire il verso. Infatti la disequazione iniziale si riduce a $-1>0$ nel dominio $x\ne0$ ed è impossibile, mentre la moltiplicazione senza cambio di verso darebbe erroneamente $x<0$.
\end{solution}
\part[1] La disequazione $(x-2)^2\ge0$ ha come insieme delle soluzioni $\mathbb{R}$.
\begin{solution}
Vero. Il quadrato è sempre non negativo ed è definito per ogni numero reale.
\end{solution}
\part[1] La disequazione $(x+1)^2+x^2\le0$ non ha soluzioni.
\begin{solution}
Vero. I due quadrati dovrebbero annullarsi contemporaneamente, ma il primo si annulla in $-1$ e il secondo in $0$.
\end{solution}
\part[1] L'insieme delle soluzioni della disequazione $\dfrac{x}{x}>0$ è $\mathbb{R}$.
\begin{solution}
Falso. Il rapporto vale $1$ soltanto nel dominio $x\ne0$. Quindi $S=\mathbb{R}\setminus\{0\}$.
\end{solution}
\end{parts}
\begin{solution}
\textit{Criteri per ciascuna risposta.} Mezzo punto per la scelta e mezzo punto per la giustificazione.
\end{solution}
\end{esercizio}
\begin{esercizio}
\obiettivo{risolvere una disequazione algebrica}
\question[4] Risolvi la disequazione.
\[\frac{(x-2)^2}{3}-\frac{(x+1)^2}{3}<2-x.\]
\begin{solution}
La differenza a sinistra vale $-2x+1$. Quindi $-2x+1<2-x$, da cui $x>-1$ e $S=(-1,+\infty)$.
\par\textit{Criteri.} Sviluppi e riduzione 2 punti, risoluzione e insieme delle soluzioni 2 punti.
\end{solution}
\end{esercizio}
\begin{esercizio}
\obiettivo{risolvere un sistema di disequazioni}
\question[4] Risolvi il sistema.
\[\begin{cases}\dfrac{x-2}{3}+x\ge2\\[5pt]\dfrac{2x+1}{4}>x-2\end{cases}\]
\begin{solution}
La prima dà $4x\ge8$, quindi $x\ge2$. La seconda dà $2x+1>4x-8$, quindi $x<\frac92$. Pertanto $S=[2,\frac92)$.
\par\textit{Criteri.} Prima disequazione 1,5 punti, seconda disequazione 1,5 punti, intersezione 1 punto.
\end{solution}
\end{esercizio}
\ifprintanswers\newpage\fi
\begin{esercizio}
\obiettivo{risolvere disequazioni algebriche}
\question Risolvi le disequazioni.
\begin{parts}
\part[5] $\displaystyle \frac{2}{x^2+x-6}-\frac{1}{x+3}\ge\frac{1}{x-2}.$
\begin{solution}
Scomponendo $x^2+x-6=(x+3)(x-2)$ e portando tutto a sinistra si ottiene
$\dfrac{1-2x}{(x+3)(x-2)}\ge0$. I valori critici sono $-3,\frac12,2$.
Da sinistra i segni del rapporto sono $+,-,+,-$. Il numeratore si annulla in $\frac12$.
Si escludono $-3$ e $2$. Quindi $S=(-\infty,-3)\cup[\frac12,2)$.
\par\textit{Criteri.} Riduzione a un'unica frazione 2 punti, studio dei segni 2 punti, dominio e insieme delle soluzioni 1 punto.
\end{solution}

\part[5] $\displaystyle (x^2-4)(x-1)>2(x+2)(x-1)^2.$
\begin{solution}
Portando a sinistra e raccogliendo $(x+2)(x-1)$, resta
$(x+2)(x-1)[x-2-2(x-1)]>0$, cioè $-x(x+2)(x-1)>0$.
Gli zeri sono $-2,0,1$. Da sinistra i segni del prodotto sono $+,-,+,-$.
Quindi $S=(-\infty,-2)\cup(0,1)$.
\par\textit{Criteri.} Riduzione alla forma fattorizzata 2 punti, studio dei segni 2 punti, insieme delle soluzioni 1 punto.
\end{solution}

\end{parts}
\end{esercizio}
\begin{esercizio}
\obiettivo{risolvere un sistema di disequazioni}
\question[6] Risolvi il sistema e rappresenta graficamente le soluzioni.
\[
\begin{cases}
1\le\dfrac9{(x-1)^2}\\[6pt]
\dfrac{1}{2x+2}<1+\dfrac{1}{x+1}
\end{cases}
\]
\begin{solution}
Per $x\ne1$ il denominatore della prima è positivo. La disequazione
equivale a $(x-1)^2\le9$, da cui $S_1=[-2,1)\cup(1,4]$.
La seconda si riduce a $\dfrac{2x+3}{2(x+1)}>0$. I punti critici sono
$-\frac32$ e $-1$, entrambi esclusi. Dunque
$S_2=(-\infty,-\frac32)\cup(-1,+\infty)$.
Ne segue $S=[-2,-\frac32)\cup(-1,1)\cup(1,4]$.
\begin{center}
\begin{tikzpicture}[x=1.23cm,y=1cm]
\draw[->,black!60] (-2.42,0)--(4.48,0) node[right] {$x$};
\draw[blue!60!black,very thick] (-2,0)--(-1.5,0);
\draw[blue!60!black,very thick] (-1,0)--(4,0);
\draw[blue!60!black,thick,fill=blue!60!black] (-2,0) circle[radius=1.7pt] node[below=3pt] {$\scriptstyle -2$};
\draw[blue!60!black,thick,fill=white] (-1.5,0) circle[radius=1.7pt] node[below=3pt] {$\scriptstyle -\frac32$};
\draw[blue!60!black,thick,fill=white] (-1,0) circle[radius=1.7pt] node[below=3pt] {$\scriptstyle -1$};
\draw[blue!60!black,thick,fill=white] (1,0) circle[radius=1.7pt] node[below=3pt] {$\scriptstyle 1$};
\draw[blue!60!black,thick,fill=blue!60!black] (4,0) circle[radius=1.7pt] node[below=3pt] {$\scriptstyle 4$};
\end{tikzpicture}
\end{center}
\par\textit{Criteri.} Prima disequazione 2 punti, seconda disequazione 2 punti, intersezione e grafico 2 punti.
\end{solution}
\end{esercizio}
\begin{esercizio}
\obiettivo{risolvere un problema geometrico}
\question[2]

Un'aiuola rettangolare ha perimetro $36$~m. La sua area deve essere almeno
$80$~m$^2$ e il lato lungo deve superare il lato corto di almeno $1$~m.
Determina quali lunghezze può avere il lato lungo.
\begin{solution}
Sia $x$ la lunghezza in metri del lato lungo. Il lato corto misura $18-x$,
con $9<x<18$. Per l'area si richiede $x(18-x)\ge80$, ossia
$(x-8)(x-10)\le0$. Quindi $8\le x\le10$.
La differenza deve soddisfare $x-(18-x)\ge1$, cioè $x\ge\frac{19}{2}$.
Il lato lungo può misurare da $9{,}5$ a $10$ metri, estremi compresi.
\par\textit{Criteri.} Modellizzazione 1 punto, risoluzione e risposta 1 punto.
\end{solution}
\end{esercizio}
\end{esercizi}
\finefila
\iniziofila{B}
\begin{esercizi}
\begin{esercizio}
\obiettivo{valutare proprietà e procedimenti}
\question Indica se ogni affermazione è vera o falsa. Giustifica la risposta.
\begin{parts}
\part[1] La disequazione $(x+2)^2+x^2\le0$ non ha soluzioni.
\begin{solution}
Vero. I due quadrati dovrebbero annullarsi contemporaneamente, ma il primo si annulla in $-2$ e il secondo in $0$.
\end{solution}
\part[1] L'insieme delle soluzioni della disequazione $\dfrac{2x}{2x}>0$ è $\mathbb{R}$.
\begin{solution}
Falso. Il rapporto vale $1$ soltanto nel dominio $x\ne0$. Quindi $S=\mathbb{R}\setminus\{0\}$.
\end{solution}
\part[1] È possibile risolvere la disequazione $\dfrac{x-2}{x}>\dfrac{2x-2}{x}$ moltiplicando ambo i membri per $x$.
\begin{solution}
Falso. Il segno di $x$ è variabile: se $x<0$ occorre invertire il verso. Nel dominio $x\ne0$ la disequazione equivale a $-1>0$, mentre la moltiplicazione senza cambio di verso darebbe erroneamente $x<0$.
\end{solution}
\part[1] La disequazione $(x-3)^2\ge0$ ha come insieme delle soluzioni $\mathbb{R}$.
\begin{solution}
Vero. Il quadrato è sempre non negativo ed è definito per ogni numero reale.
\end{solution}
\end{parts}
\begin{solution}
\textit{Criteri per ciascuna risposta.} Mezzo punto per la scelta e mezzo punto per la giustificazione.
\end{solution}
\end{esercizio}
\begin{esercizio}
\obiettivo{risolvere una disequazione algebrica}
\question[4] Risolvi la disequazione.
\[\frac{(x+3)^2}{3}-\frac{(x-1)^2}{3}\ge x+2.\]
\begin{solution}
La differenza a sinistra vale $\frac{8x+8}{3}$. Quindi $8x+8\ge3x+6$, da cui $x\ge-\frac25$ e $S=[-\frac25,+\infty)$.
\par\textit{Criteri.} Sviluppi e riduzione 2 punti, risoluzione e insieme delle soluzioni 2 punti.
\end{solution}
\end{esercizio}
\begin{esercizio}
\obiettivo{risolvere un sistema di disequazioni}
\question[4] Risolvi il sistema.
\[\begin{cases}\dfrac{x+2}{3}+x\le-2\\[5pt]\dfrac{-2x+1}{4}>-x-2\end{cases}\]
\begin{solution}
La prima dà $4x\le-8$, quindi $x\le-2$. La seconda dà $-2x+1>-4x-8$, quindi $x>-\frac92$. Pertanto $S=(-\frac92,-2]$.
\par\textit{Criteri.} Prima disequazione 1,5 punti, seconda disequazione 1,5 punti, intersezione 1 punto.
\end{solution}
\end{esercizio}
\ifprintanswers\newpage\fi
\begin{esercizio}
\obiettivo{risolvere disequazioni algebriche}
\question Risolvi le disequazioni.
\begin{parts}
\part[5] $\displaystyle \frac{2}{x^2+3x-4}-\frac1{x+4}\ge\frac1{x-1}.$
\begin{solution}
Scomponendo $x^2+3x-4=(x+4)(x-1)$ e portando tutto a sinistra si ottiene
$\dfrac{-1-2x}{(x+4)(x-1)}\ge0$. I valori critici sono $-4,-\frac12,1$.
Da sinistra i segni del rapporto sono $+,-,+,-$. Il numeratore si annulla in $-\frac12$.
Si escludono $-4$ e $1$. Quindi $S=(-\infty,-4)\cup[-\frac12,1)$.
\par\textit{Criteri.} Riduzione a un'unica frazione 2 punti, studio dei segni 2 punti, dominio e insieme delle soluzioni 1 punto.
\end{solution}

\part[5] $\displaystyle (x^2-9)(x-2)>2(x+3)(x-2)^2.$
\begin{solution}
Portando a sinistra e raccogliendo $(x+3)(x-2)$, resta
$(x+3)(x-2)[x-3-2(x-2)]>0$, cioè $(x+3)(x-2)(1-x)>0$.
Gli zeri sono $-3,1,2$. Da sinistra i segni del prodotto sono $+,-,+,-$.
Quindi $S=(-\infty,-3)\cup(1,2)$.
\par\textit{Criteri.} Riduzione alla forma fattorizzata 2 punti, studio dei segni 2 punti, insieme delle soluzioni 1 punto.
\end{solution}

\end{parts}
\end{esercizio}
\begin{esercizio}
\obiettivo{risolvere un sistema di disequazioni}
\question[6] Risolvi il sistema e rappresenta graficamente le soluzioni.
\[
\begin{cases}
1\le\dfrac{16}{(x+1)^2}\\[6pt]
\dfrac{1}{2x-2}<1+\dfrac{1}{x-1}
\end{cases}
\]
\begin{solution}
Per $x\ne-1$ il denominatore della prima è positivo. La disequazione
equivale a $(x+1)^2\le16$, da cui $S_1=[-5,-1)\cup(-1,3]$.
La seconda si riduce a $\dfrac{2x-1}{2(x-1)}>0$. I punti critici sono
$\frac12$ e $1$, entrambi esclusi. Dunque
$S_2=(-\infty,\frac12)\cup(1,+\infty)$.
Ne segue $S=[-5,-1)\cup(-1,\frac12)\cup(1,3]$.
\begin{center}
\begin{tikzpicture}[x=1.23cm,y=1cm]
\draw[->,black!60] (-5.43,0)--(3.49,0) node[right] {$x$};
\draw[blue!60!black,very thick] (-5,0)--(0.5,0);
\draw[blue!60!black,very thick] (1,0)--(3,0);
\draw[blue!60!black,thick,fill=blue!60!black] (-5,0) circle[radius=1.7pt] node[below=3pt] {$\scriptstyle -5$};
\draw[blue!60!black,thick,fill=white] (-1,0) circle[radius=1.7pt] node[below=3pt] {$\scriptstyle -1$};
\draw[blue!60!black,thick,fill=white] (0.5,0) circle[radius=1.7pt] node[below=3pt] {$\scriptstyle \frac12$};
\draw[blue!60!black,thick,fill=white] (1,0) circle[radius=1.7pt] node[below=3pt] {$\scriptstyle 1$};
\draw[blue!60!black,thick,fill=blue!60!black] (3,0) circle[radius=1.7pt] node[below=3pt] {$\scriptstyle 3$};
\end{tikzpicture}
\end{center}
\par\textit{Criteri.} Prima disequazione 2 punti, seconda disequazione 2 punti, intersezione e grafico 2 punti.
\end{solution}
\end{esercizio}
\begin{esercizio}
\obiettivo{risolvere un problema geometrico}
\question[2]

Un'aiuola rettangolare ha perimetro $44$~m. La sua area deve essere almeno
$120$~m$^2$ e il lato lungo deve superare il lato corto di almeno $1$~m.
Determina quali lunghezze può avere il lato lungo.
\begin{solution}
Sia $x$ la lunghezza in metri del lato lungo. Il lato corto misura $22-x$,
con $11<x<22$. Per l'area si richiede $x(22-x)\ge120$, ossia
$(x-10)(x-12)\le0$. Quindi $10\le x\le12$.
La differenza deve soddisfare $x-(22-x)\ge1$, cioè $x\ge\frac{23}{2}$.
Il lato lungo può misurare da $11{,}5$ a $12$ metri, estremi compresi.
\par\textit{Criteri.} Modellizzazione 1 punto, risoluzione e risposta 1 punto.
\end{solution}
\end{esercizio}
\end{esercizi}
\finefila
\end{document}
